\metatitle{Real-rooted polynomials: tableaux and polytopes}

\metadescription{Real-rootedness for descent polynomials of standard Young tableaux, Ehrhart polynomials from Schur functions, and related combinatorial families.}


\section[realRootedTableaux]{Tableaux and polytopes}

This page collects real-rootedness results from tableaux and polytope
combinatorics: descent polynomials of standard Young tableaux, Ehrhart
polynomials from Schur functions, and related families.  The main tools are
\hyperref[interlacingPolynomials]{interlacing},
\hyperref[brentiEhrhartCriterion]{Brenti's Ehrhart criterion}, and occasionally
\hyperref[polyaFrequency]{total-positivity} methods.  For the broader context,
see \hyperref[polynomialRootIntro]{unimodality and real-rootedness}.

\subsection[brentiEhrhartCriterion]{Brenti's Ehrhart criterion}

The examples below use the same bridge between Ehrhart theory and
real-rootedness.  Suppose \(L(m)\) is a polynomial of degree \(d\), and define
the numerator \(H(t)\) by
\[
  \sum_{m\geq 0} L(m)t^m = \frac{H(t)}{(1-t)^{d+1}}.
\]
In the factored cases below, the hypotheses in Brenti's
Worpitzky-transform criterion \cite[Thm.~4.6.1]{Brenti1989} are checked
directly from the linear-factor product.  The conclusion is that the numerator
\(H(t)\) has only real roots.  Thus the proof strategy is to identify the
desired descent or \(h^*\)-polynomial as such a numerator and then prove that
\(L(m)\) factors with only real roots.

\subsection[sytDescentRealRooted]{Descent polynomials of standard Young tableaux}

\begin{example*}[Descents in standard Young tableaux]

For a non-skew shape $\lambda$ with $n$ boxes, the descent-generating polynomial
($i$ is a descent if $i+1$ appears in a strictly lower row)
\[
   W_\lambda(t) \coloneqq \sum_{T \in \SYT(\lambda)} t^{\des(T)}
\]
is real-rooted \cite[Thm. 5.2.3]{Brenti1989}.
The argument proceeds in two steps.
First, \cite[Eq. (7.96)]{StanleyEC2} gives the identity
\[
  \sum_{m \geq 0} \schurS_{\lambda}(1^m)\,z^m = \frac{W_\lambda(z)}{(1-z)^{n+1}},
\]
which identifies $W_\lambda(z)$ as the \hyperref[ehrhart]{$h^*$-polynomial} of the
associated order polytope, and $m \mapsto \schurS_\lambda(1^m)$ as its Ehrhart polynomial.
Second, the hook-content formula gives
\[
  \schurS_\lambda(1^m) = \prod_{u \in \lambda} \frac{m + c(u)}{h(u)},
\]
where $c(u)$ is the content and $h(u)$ the hook length of cell $u$.
Each factor is linear in $m$, so the relevant Ehrhart/order polynomial is
real-rooted.  By
\hyperref[brentiEhrhartCriterion]{Brenti's Ehrhart criterion}, the numerator
is real-rooted, giving real-rootedness of $W_\lambda(t)$.

This argument does not extend to skew shapes: the Ehrhart polynomial of $\lambda/\mu$
need not be real-rooted, so Brenti's criterion does not apply (see the remark below).

\cite{Branden2004operators} shows moreover that $W_\lambda(t)$ \hyperref[interlacingPolynomials]{interlaces}
$W_{\lambda^+}(t)$ whenever $\lambda^+$ is obtained from $\lambda$ by adding a box.
In particular, for $\lambda=(n,n)$, $W_\lambda(t)$ is the \hyperref[realRootedCatalan]{Narayana polynomial} $N_n(t)$.
\end{example*}


\begin{remark*}[Skew shapes]
For skew shapes $\lambda/\mu$, the identity $\sum_{m\geq 0} \schurS_{\lambda/\mu}(1^m)\,z^m = W_{\lambda/\mu}(z)/(1-z)^{n+1}$
still holds, but the Ehrhart polynomial $m \mapsto \schurS_{\lambda/\mu}(1^m)$ need not be real-rooted,
so \hyperref[brentiEhrhartCriterion]{Brenti's criterion} does not apply.
Whether $W_{\lambda/\mu}(t)$ is nevertheless real-rooted for all skew shapes is open.
\end{remark*}

\subsection[schurEhrhartRealRooted]{Schur Ehrhart polynomials}

\begin{example*}[Ehrhart polynomials from Schur functions]

Let $k = \ell(\lambda)$ be the number of parts of $\lambda$.
The map
\[
 n \mapsto |\SSYT(n \lambda, k)|,
\]
counting semi-standard Young tableaux of shape $n\lambda$ with entries in $\{1,\dotsc,k\}$,
is a polynomial in $n$. The Weyl character formula gives
\[
  P_\lambda(n) \coloneqq  \prod_{1 \leq i \lt j \leq k} \frac{n(\lambda_i - \lambda_j)+j-i}{j-i}.
\]
Since $\lambda_i \geq \lambda_j$ for $i < j$, each factor is linear in $n$ with a non-positive root,
so $P_\lambda(n)$ is real-rooted as a product of such factors.
By \hyperref[brentiEhrhartCriterion]{Brenti's Ehrhart criterion}, this implies
that the corresponding $h^*$-polynomial is also real-rooted.
Brändén \cite{Branden2004operators} showed moreover that $P_\lambda(n)$
\hyperref[interlacingPolynomials]{interlaces} $P_\mu(n)$ whenever $\mu$ covers $\lambda$ in Young's lattice.
The above map is the Ehrhart polynomial of the \hyperref[gtpatterns]{Gelfand--Tsetlin polytope} associated with $\lambda$ and $k$.


We note that the map
\[
 n \mapsto |\SSYT(n \lambda/ n\mu, k)|
\]
for skew shapes is not real-rooted in general: $\lambda =(3,2)$, $\mu = (1)$ gives for $k=3$
the polynomial $\frac{1}{4} (n+1)^2 \left(7 n^2+10 n+4\right)$ which has non-real roots.

It does look like the corresponding $h^*$-polynomials are real-rooted even for skew shapes.
For non-skew shapes, this is proved by \name{Francesco Brenti} \cite{Brenti1989}.
\end{example*}
